\section{Overview and Motivation (Numerical Methods)}\label{sec:c5s1}
\lipsum[138-138]

This section builds on the material in \cref{ch:5}, \cref{ch:6} and the prior work of \citet{ch5-4}; see also \citep{ch5-14,ch5-8}.

\begin{equation}\label{eq:c5s1-a}
  I = \int_{0}^{1}\!\!\int_{0}^{1} \frac{\sin(\pi x)\sin(\pi y)}{1+x^2+y^2}\,\mathrm{d}x\,\mathrm{d}y \approx \frac{1}{M}\sum_{m=1}^{M} h(\xi_m),
\end{equation}
with variance decreasing as
\begin{equation}\label{eq:c5s1-b}
  \operatorname{Var}\bigl[\hat I_M\bigr] = \frac{1}{M}\Bigl(\int h^2\,\mathrm{d}\mu - I^2\Bigr) = \mathcal{O}(M^{-1}).
\end{equation}

Equations~\eqref{eq:c5s1-a} and~\eqref{eq:c5s1-b} are used throughout the remainder of this section.

\kant[2-10]

\subsection{Quantitative behaviour}\label{sec:c5s1-q}
Results are collected in \cref{tab:c5s1} and visualised in \cref{fig:c5s1-f1}. Similar trends were reported in~\citep{ch5-6,ch5-10}.

\begin{table}[htbp]
\centering
\caption{Summary of measured quantities for experiment set \texttt{c5s1}.}\label{tab:c5s1}
\begin{tabular}{lrrr}
\toprule
Method & Score & Latency & Error \\
\midrule
Alpha & \num{70.9} & \SI{5.61}{\milli\second} & \num{4.447e-5} \\
Beta & \num{70.9} & \SI{4.53}{\milli\second} & \num{8.386e-2} \\
Gamma & \num{67.1} & \SI{8.61}{\milli\second} & \num{2.016e-1} \\
Delta & \num{36.8} & \SI{9.37}{\milli\second} & \num{1.091e-5} \\
Epsilon & \num{15.7} & \SI{5.01}{\milli\second} & \num{8.377e-5} \\
\bottomrule
\end{tabular}
\end{table}

\begin{figure}[htbp]
\centering
\begin{tikzpicture}[node distance=9mm and 8mm,>=Stealth,blk/.style={draw,rounded corners,fill=orange!12,minimum width=2.1cm,minimum height=8mm,align=center}]
\node[blk](a){Data\\acquisition};\node[blk,right=of a](b){Pre-\\processing};\node[blk,right=of b](c){Model\\fitting};\node[blk,right=of c](d){Validation};
\node[blk,below=of b](e){Feature\\store};\node[blk,below=of c](f){Parameter\\search};
\draw[->](a)--(b);\draw[->](b)--(c);\draw[->](c)--(d);\draw[->](b)--(e);\draw[->](e)--(f);\draw[->](f)--(c);\draw[->,dashed](d.north)--++(0,.8)-|(a.north);
\begin{scope}[on background layer]\node[draw,dashed,fit=(a)(b)(c)(d)(e)(f),inner sep=6pt,fill=gray!5]{};\end{scope}
\end{tikzpicture}
\caption{Generated figure for \texttt{c5s1-f1} (parameters $a=3$, $b=3$).}\label{fig:c5s1-f1}
\end{figure}

\kant[3-10]

\subsection{Implementation notes}\label{sec:c5s1-i}
The core routine is shown in \cref{lst:c5s1}; its complexity is $\mathcal{O}(N\log N)$ per iteration~\cite{ch5-13}.

\begin{lstlisting}[language=python,caption={Reference implementation for c5s1.},label={lst:c5s1}]
def conjugate_gradient(A, b, tol=1e-8, maxit=500):
    x = b * 0
    r = b - A @ x
    p = r.copy()
    for k in range(maxit):
        Ap = A @ p
        alpha = (r @ r) / (p @ Ap)
        x += alpha * p
        r_new = r - alpha * Ap
        if (r_new @ r_new) ** 0.5 < tol:
            break
        p = r_new + ((r_new @ r_new) / (r @ r)) * p
        r = r_new
    return x
\end{lstlisting}

\lipsum[68-68]
\begin{theorem}\label{thm:c5s1}
Let $A\in\mathbb{R}^{n\times n}$ be symmetric positive definite with condition number $\kappa$. Then the iteration converges and
$\lVert e_k\rVert_A \le 2\bigl(\tfrac{\sqrt\kappa-1}{\sqrt\kappa+1}\bigr)^k\lVert e_0\rVert_A$.
\end{theorem}
\begin{enumerate}[label=(\roman*)]
\item the bound in \cref{thm:c5s1} is sharp for some initial errors;
\item the constant $2$ cannot be improved in general;
\item preconditioning reduces the effective $\kappa$.
\end{enumerate}

\begin{figure}[htbp]
\centering
\includegraphics[width=0.45\linewidth]{diag.png}
\caption{Raster/vector asset \texttt{diag.png} included from the figures directory.}\label{fig:c5s1-i0}
\end{figure}

\kant[6-11]
